\chapter{Analyse des jeux de données}
\label{chap:datasets}

Le choix des données conditionne la pertinence d'une solution de maintenance prédictive.
Plutôt que de se limiter à un seul cas, le projet s'appuie sur \textbf{cinq jeux de données
de référence}, couvrant des natures de machines, des modalités d'acquisition et des tâches
différentes. Cette diversité valide le caractère \emph{générique} du pipeline.

\section{Vue comparative}
Le tableau~\ref{tab:datasets} synthétise les caractéristiques techniques réelles des
datasets après traitement (dimensions effectives mesurées sur les tenseurs produits).

\begin{table}[h]
\centering
\small
\caption{Caractéristiques des jeux de données utilisés (valeurs réelles après pipeline).}
\label{tab:datasets}
\begin{tabular}{>{\bfseries}l c c c c c}
\toprule
Dataset & Tâche & Fréq. éch. & Canaux & Fenêtres & Features \\
\midrule
VBL-VA001         & Classif. (4 cl.)  & 20\,kHz   & 3 (x,y,z) & 400    & 45  \\
CWRU              & Classif. (10 cl.) & 12\,kHz   & 1 (DE)    & 12\,627 & 15  \\
Mechanical Faults & Classif. (4 cl.)  & 25\,kHz   & 4         & 36\,000 & 132 \\
CMAPSS FD001      & Régression (RUL)  & cycles    & 21 capt.  & 17\,731 & 126 \\
MCC5-THU          & Classif. (8 cl.)  & 12{,}8\,kHz & 3 (x,y,z) & ---    & ---  \\
\bottomrule
\end{tabular}
\end{table}

\section{VBL-VA001 --- bancs d'essai vibratoires}
Ce jeu de données fournit des signaux tri-axiaux (x, y, z) échantillonnés à 20\,kHz sur un
banc d'essai, répartis en \textbf{quatre classes parfaitement équilibrées} (100 fenêtres
chacune) : \texttt{normal}, \texttt{bearing} (roulement), \texttt{misalignment}
(désalignement) et \texttt{unbalance} (balourd). Après extraction, chaque fenêtre est
décrite par \textbf{45 features} (15 par canal). Sa petite taille (400 fenêtres) et la
séparabilité nette de ses classes en font un cas \emph{facile}, idéal pour valider le
pipeline de bout en bout.

\section{CWRU --- défauts de roulement (Case Western Reserve University)}
Référence académique mondiale du diagnostic de roulement, CWRU fournit des signaux
d'accélération (canal \texttt{DE\_time}, côté entraînement) échantillonnés à 12\,kHz. Le
projet en retient une taxonomie fine à \textbf{10 classes} croisant la \emph{localisation}
du défaut (bague interne, bague externe, bille) et sa \emph{sévérité} (faible, moyen,
grave), plus l'état sain :
\begin{center}\small\ttfamily
sain, roulement\_interne\_\{faible,moyen,grave\}, roulement\_externe\_\{faible,moyen,grave\},
roulement\_bille\_\{faible,moyen,grave\}
\end{center}
Le dataset traité compte \textbf{12\,627 fenêtres} et \textbf{15 features} (un seul canal).
La classe \texttt{sain} domine (3\,531 fenêtres) tandis que les neuf classes de défaut
comptent environ 1\,010 fenêtres chacune. Cette granularité (localisation + sévérité)
rend la tâche réaliste et exigeante.

\section{Mechanical Faults --- machine tournante multi-capteurs}
Ce jeu de données provient d'un banc de machine tournante instrumenté de \textbf{4
accéléromètres} (vertical et horizontal, côté poulie et côté disque), échantillonnés à
25\,kHz. Quatre classes équilibrées de \textbf{9\,000 fenêtres} chacune (36\,000 au total)
sont considérées : \texttt{sain}, \texttt{desalignement}, \texttt{desequilibre} et
\texttt{jeu\_mecanique}. La richesse multi-capteurs et la finesse spectrale conduisent à un
vecteur de \textbf{132 features} (33 par canal, dont des descripteurs harmoniques et
d'enveloppe spécialisés). C'est le dataset le plus volumineux et le plus difficile pour la
classification (cf. chapitre~\ref{chap:bench}).

\section{CMAPSS FD001 --- pronostic de durée de vie (NASA)}
CMAPSS (\emph{Commercial Modular Aero-Propulsion System Simulation}) est la référence du
\textbf{pronostic de RUL}. Le sous-ensemble FD001 simule la dégradation de \textbf{100
moteurs} jusqu'à la panne, chacun suivi par \textbf{21 capteurs} au fil des cycles de
fonctionnement. Le projet en dérive \textbf{17\,731 fenêtres} décrites par \textbf{126
features} (statistiques sur fenêtres de cycles des 21 capteurs). La cible principale est la
RUL (durée de vie restante, plafonnée à 125 cycles, moyenne $\approx 80{,}6$). Des labels de
dégradation (\texttt{sain}, \texttt{degradation\_precoce}, \texttt{degradation\_avancee},
\texttt{critique}) sont également dérivés pour la détection d'anomalies. À noter : les
données brutes comportaient 119\,362 valeurs manquantes (capteurs constants), traitées par
remplacement à zéro lors du benchmark.

\section{MCC5-THU --- défauts d'engrenage}
Ce jeu de données de réducteur (gearbox), échantillonné à 12{,}8\,kHz sur 3 axes, comporte
\textbf{huit classes} de défauts d'engrenage : sain, \texttt{pitting} (piqûres),
\texttt{usure}, \texttt{dent\_cassee}, \texttt{fissure}, \texttt{dent\_manquante} et deux
défauts \emph{mixtes} (rupture de dent combinée à un défaut de roulement interne/externe).
Il est traité par un \textbf{réseau de neurones convolutif 1D} (CNN) opérant directement sur
le signal, sans extraction manuelle de features (cf. chapitre~\ref{chap:bench}).

\section{Pipeline d'unification}
La diversité des formats (CSV, MAT, TXT, multi-canaux, fréquences hétérogènes) est absorbée
par un \textbf{adaptateur unifié} (\texttt{DatasetAdapter}) qui produit une structure
commune \texttt{UnifiedSignal}. Les opérations clés sont : le \textbf{rééchantillonnage} à
une fréquence cible de 12\,800\,Hz (\texttt{resample\_poly}), la \textbf{normalisation}
Z-score par canal, et le \textbf{fenêtrage glissant} (fenêtres de 1\,024 points, recouvrement
de 50\,\%). Un mécanisme de cache (clé MD5 sur chemin + configuration) évite de recalculer
les features. Cette abstraction est ce qui permet au même pipeline de traiter indifféremment
un roulement CWRU et un turbomoteur CMAPSS.

\section{Couverture des modes de défaillance : un spectre complet}
\label{sec:couverture}
Le véritable apport du choix des données réside dans leur \textbf{complémentarité}. Là où
la plupart des travaux académiques se concentrent sur un \emph{unique} dataset et une
\emph{unique} tâche (typiquement la classification CWRU), PrognoSense couvre l'ensemble du
spectre des défaillances des machines tournantes industrielles. Le
tableau~\ref{tab:fautes} recense les \textbf{modes de défaillance réellement intégrés} et
diagnosticables par la plateforme.

\begin{table}[h]
\centering\small
\caption{Catalogue des modes de défaillance couverts par PrognoSense.}
\label{tab:fautes}
\begin{tabular}{>{\bfseries}p{3.4cm} p{6.4cm} p{3.6cm}}
\toprule
Famille de défaut & Modes couverts & Source / signature \\
\midrule
Balourd (déséquilibre) & déséquilibre rotor & VBL, MF --- pic $1\times$ \\
Désalignement & parallèle / angulaire & VBL, MF --- pics $2\times$, $3\times$ \\
Jeu mécanique & desserrage, jeu de palier & MF --- harmoniques multiples \\
Roulement --- bague externe & faible / moyen / grave & CWRU --- BPFO + enveloppe \\
Roulement --- bague interne & faible / moyen / grave & CWRU --- BPFI + bandes latérales \\
Roulement --- bille & faible / moyen / grave & CWRU --- BSF + FTF \\
Engrenage --- piqûres & \texttt{pitting} & MCC5 --- GMF + sidebands \\
Engrenage --- usure & usure de denture & MCC5 \\
Engrenage --- dent & dent cassée, fissurée, manquante & MCC5 \\
Défauts mixtes & rupture de dent + roulement int./ext. & MCC5 \\
Dégradation turbomoteur & précoce / avancée / critique + RUL & CMAPSS (21 capteurs) \\
Anomalie inconnue & tout écart au régime sain & auto-encodeur (tous datasets) \\
\bottomrule
\end{tabular}
\end{table}

Cette couverture appelle plusieurs observations qui fondent le caractère distinctif du
projet :
\begin{itemize}[leftmargin=1.5em]
  \item \textbf{Multi-domaine} : roulements (CWRU), engrenages (MCC5), machines tournantes
        génériques (VBL, MF) et turbomoteurs (CMAPSS) sont traités par une seule plateforme.
  \item \textbf{Multi-tâche} : classification de défaut, classification de \emph{sévérité}
        (CWRU : faible/moyen/grave), régression de durée de vie (CMAPSS) et détection
        d'anomalie non supervisée coexistent.
  \item \textbf{Granularité de sévérité} : peu de travaux distinguent l'intensité d'un
        même défaut ; CWRU est ici exploité à 10 classes croisant localisation et gravité.
  \item \textbf{Défauts combinés} : MCC5 inclut des défauts \emph{mixtes} (denture +
        roulement), bien plus proches de la réalité industrielle que les défauts isolés.
  \item \textbf{Détection de l'inconnu} : au-delà des défauts catalogués, l'auto-encodeur
        repère toute déviation jamais vue, là où un classifieur supervisé reste aveugle aux
        modes absents de son entraînement.
\end{itemize}
Cette amplitude --- rare dans un projet de fin d'études --- démontre que l'architecture
n'est pas un démonstrateur mono-cas, mais une \textbf{plateforme générique} apte à
absorber de nouveaux équipements et de nouveaux défauts, ce qui constitue son principal
facteur de différenciation technique.
